\section{Conclusiones y Trabajo a Futuro}%
\label{sec:conclusions}

La organización eficiente de la información vinculada a viajes suele ser un desafío, debido a que suelen estar
involucrados elementos (vuelos, trenes, hoteles, etc.) de diversos proveedores distintos que, además, suele recibirse a
través de correos electrónicos con diferentes interfaces.

Como consecuencia, el usuario, al momento de querer consultar el itinerario de un viaje, debe consultar cada uno de esos
mails individualmente o interactuar con las respectivas aplicaciones de cada proveedor, dificultando una vista integral
del itinerario completo.
Con este panorama, en muchos casos, el usuario termina aplicando como solución la copia de toda esa información en otro
lugar o documento, con el riesgo de cometer un error o de contar con información desactualizada.

En este contexto, la propuesta de este trabajo fue ofrecer una aplicación capaz de extraer información relevante de
correos electrónicos de reservas, independientemente de la estructura o el remitente del correo, para luego proceder a
la creación de un itinerario integral del viaje.

El presente proyecto se centrará en la extracción de información relevante contenida exclusivamente en correos
electrónicos de reservas de tickets aéreos y tickets de trenes, que es extensible a otros dominios, como estadías de
hotel, alquiler de vehículos, entradas para espectáculos, excusiones, etc.

El objetivo principal fue implementar un modelo que permita generar un itinerario integral por usuario, centralizando la
información relevante de sus viajes mediante la extracción, identificación y clasificación automática de datos
contenidos en correos electrónicos de confirmación (como reservas de vuelos, alojamientos, excursiones, entre otros),
independientemente del proveedor del servicio y de la estructura del mensaje.
Adicionalmente, se le agregó como funcionalidad complementaria la generación de recomendaciones de puntos de interés en
los destinos identificados, con el objetivo de enriquecer la experiencia del usuario y aportar valor adicional al
itinerario construido.

La presente implementación comienza con la extracción automática de los campos relevantes de correos de vuelos y trenes,
independientemente de la estructura HTML, es decir, solo analizando texto plano sin utilizar marcas estructurales que
actúen como indicios (como atributos de nodos HTML, asunto del correo electrónico, etc.).

Con el uso de CoreNLP se pudo ver que el proceso de Named Entity Recognition es solo el primer paso para una correcta
identificación.
En el caso de la identificación de segmentos de vuelos o trenes, fue necesario ubicar en el texto a la entidad
reconocida y, en función de la proximidad con otras entidades reconocidas, decidir con cierto nivel de confianza la
información a extraer.

El prototipo funcional desarrollado permitió validar las técnicas aplicadas de extracción de información, a través de
una solución web simple que simuló ser la vista del usuario de la aplicación de gestión de itinerarios, con un grado de
usabilidad apropiado.

En cuanto al trabajo futuro, una primera mejora posible consiste en la implementación de un clasificador de correos
electrónicos, similar a como se implementa un filtro de spam.
El clasificador produciría un valor de confianza para decidir en tiempo real si el correo corresponde a una reserva de
vuelos, de trenes o de otro tipo, y así optar por el pipeline apropiado cuando el usuario realiza drag \& drop de
cualquier correo en una sola área.

También se podrían incorporar otros pipelines que cubran dominios no alcanzados por la implementación actual, como
estadías en hoteles, alquileres de autos, atracciones y reservas en restaurantes.
También se puede examinar la implementación de un pipeline en español para los dominios de vuelos y trenes.
Para cada uno de estos nuevos pipelines se deberán abordar una serie de desafíos propios de cada dominio concreto.

Finalmente, respecto a la usabilidad, una mejora técnica posible es sustituir la tarea de exportar un correo a .eml para
que, en cambio, el usuario reenvíe los correos a una casilla como \texttt{trips@email2trip.xyz}, de forma que la
aplicación analice los correos recibidos y los asigne según el remitente, que debe estar registrado en la aplicación con
esa misma dirección de correo.
